% Source: https://artofproblemsolving.com/community/c3h2580730_1998_hexagons_666_dont_share_a_vertex_ukraine_1998_grade_8_p2
% Stable identity: aops-topic-2580730-q01
In the hexagonal tiling, two hexagons that share a vertex also share an edge. Colour the hexagons with three colours so that neighbouring hexagons get different colours: in axial coordinates $(q,s)$ (the six neighbours differ by $(\pm1,0)$, $(0,\pm1)$, $\pm(1,-1)$), give the hexagon $(q,s)$ the colour $(q-s)\bmod3$. The difference $q-s$ changes by $\pm1$ or $\pm2$ between neighbours, never by a multiple of $3$, so neighbours get different colours.
\begin{asy}
usepackage("amsmath");
size(5cm);
defaultpen(fontsize(10pt));
real r=0.55;pen[] col={rgb(0.85,0.3,0.3),rgb(0.3,0.6,0.85),rgb(0.95,0.8,0.3)};
for(int q=-3;q<=3;++q)for(int s=-2;s<=2;++s){pair c=(sqrt(3)*r*(q+s/2),1.5*r*s);
if(abs(c.x)>2.2) continue;
path h=c+r*dir(30)--c+r*dir(90)--c+r*dir(150)--c+r*dir(210)--c+r*dir(270)--c+r*dir(330)--cycle;
int k=((q-s)%3+3)%3;filldraw(h,col[k],black);label(string(k+1),c);}
\end{asy}

The $1998$ marked hexagons fall into $3$ colour classes, so by the pigeonhole principle some class contains at least $\dfrac{1998}{3}=666$ marked hexagons. Choose $666$ of them. They all have the same colour, so no two of them are neighbours, and no two share a vertex. $\blacksquare$
